\subsection{在连续映射下的像}

命题 4. --- 设 E 与 F 是两个局部凸空间，\(f:E \rightarrow F\) 是一个连续映射。假设 \(f(0) = 0\) ，且存在一个实数 \(m \geqslant 0\) ，使得 \(f(\lambda x) = \lambda^{m}f(x)\) 对每个 \(\lambda > 0\) 成立。那么，如果 A 是 E 的一个有界子集，则 \(f(A)\) 在 F 中有界。

事实上，如果 V 是 F 中 0 的一个邻域，那么 \(f^{-1}(V)\) 是 E 中 0 的一个邻域。如果 A 在 E 中有界，则存在 \(\lambda > 0\) ，使得 \(A \subset \lambda f^{-1}(V)\) ，这蕴含 \(f(A) \subset \lambda^{m}V\) 。

推论 1. --- 设 E 与 F 是两个局部凸空间，\(u:E\to F\) 是一个连续线性映射。如果 A 是 E 的一个有界子集，那么 \(u(A)\) 在 F 中有界。

推论 2. --- 设 \(\mathrm{E} = \prod_{i\in \mathrm{I}}\mathrm{E}_i\) 是一个局部凸空间族的乘积。

那么 E 的一个子集有界，当且仅当它的所有投影都有界。更一般地：

推论 3. --- 设 E 是一个向量空间，\((\mathrm{F}_{i})_{i\in\mathrm{I}}\) 是一个局部凸空间族，\(f_{i}\) 是从 E 到 \(F_{i}\) 中的一个线性映射（对每个 \(i\in I\) ）。假设给 E 装备使所有 \(f_{i}\) 都连续的最粗（局部凸）拓扑（II, 第 26 页）。那么，要使 E 的子集 A 有界，必须且只需 \(f_{i}(A)\) 在 \(F_{i}\) 中有界（对所有 \(i\in I\) ）。

事实上，如果 A 有界，那么诸 \(f_{i}(\mathbf{A})\) 也有界（推论 1）。反之，如果诸 \(f_{i}(\mathbf{A})\) 有界，且 p 是 E 上的一个连续半范数，那么存在 I 的一个有限子集 J 与一个族 \((q_{j})_{j\in \mathbb{J}}\) ，其中 \(q_{j}\) 是 \(\mathbf{F}_j\) 上的一个连续半范数，使得 \(p\leqslant \sup_{j\in \mathbb{J}}(q_j\circ f_j)\) ，因而 \(p\) 在 \(\mathbf{A}\) 上有界。

推论 4. --- 设 \(\mathbf{E}_i\) ( \(1 \leqslant i \leqslant n\) ) 与 \(\mathbf{F}\) 是局部凸空间，\(f\) 是从 \(\prod_{i=1}^{n} \mathbf{E}_i\) 到 \(\mathbf{F}\) 中的一个连续多重线性映射。如果 \(\mathbf{B}_i\) 是 \(\mathbf{E}_i\) 的一个有界子集（对 \(1 \leqslant i \leqslant n\) ），那么 \(f\left(\prod_{i=1}^{n} \mathbf{B}_i\right)\) 在 \(\mathbf{F}\) 中有界。

推论 5. --- 设 E 与 F 是两个局部凸空间，\(u:E\to F\) 是一个连续多项式映射。如果 A 是 E 的一个有界子集，那么 \(u(A)\) 有界。

